\section{The decision task and the data}\label{sec:data}

\subsection{The tasks and the example format}
% source: task definitions DCM:86-92; system message at index 0, DC:161; domains and formats RP:63-66.
\kleos{} routes its judgment calls through seven decision tasks: notification prioritization (``rank competing items by urgency, evidence and impact''), tool routing (``choose the right tool, or none''), mission-control briefing (``synthesize a structured briefing from context''), context prioritization (``select and order the most relevant context''), recommendation generation (``recommend actions grounded in supplied evidence''), memory-conflict resolution (``resolve contradictions using recency and reliability'') and workspace reasoning (``reason within workspace scoping rules''). An example is a conversation: a system message at index 0 that states the model's role and the rule it must apply; a user message that presents the items in one of four formats, bullets, prose, JSON or a Slack-style thread; and the reference assistant answer. The answer either commits to a ranking, with a one-line rationale per item, or declines and says what it needs; the reference carries the label of the factor that decided it, and in release v0.0.7 every training and validation answer also names that factor in a closing line. The 1,350 examples in the releases used here span four domains (career 335, coursework 341, projects 340, research 334) and the four formats (bullets 473, prose 464, JSON 349, Slack-style 64). Appendix~\ref{app:format} shows the format on a public development fixture.

\subsection{The policy}
% source: formula and constants SC:182-192, 209-215; weights SC:239-252; difficulty SC:269-271; policies SC:196-207.
Each item is scored as
\begin{equation}
\text{score} = \text{deadline\_score} \times \text{evidence\_weight} \times \text{impact\_weight}, \qquad \text{deadline\_score} = \frac{1}{1 + \text{days to the deadline}},
\end{equation}
with evidence weights \code{confirmed} 1.0, \code{corroborated} 0.8, \code{reported} 0.55, \code{single\_source} 0.35, \code{unverified} 0.2 and \code{contradicted} 0.1, and impact weights \code{high} 1.0, \code{medium} 0.6 and \code{low} 0.3. Ties are broken by the item key, so every ranking is total and reproducible. Twelve named policies combine the score with checks; Table~\ref{tab:policies} in Appendix~\ref{app:policies} lists each with what it orders by, when it declines and the deciding factor it names. Three named constants govern the checks: a close call is a relative separation $|a-b|/\max(a,b)$ below 0.15 between the top two scores; an explicit statement is stale after 30 days; and a record that challenges an explicit statement must carry evidence of at least 0.8 (\code{corroborated}). ``Adequate'' evidence means a weight of at least 0.55 (\code{reported}), and an ``explicit statement'' is an item graded \code{confirmed}. An example's difficulty is computed from the policy's own ordering rather than copied from an axis. It depends on the first sort component on which the top two candidates differ: a gating criterion such as scope or adequacy makes the example easy; the primary ranking component makes it easy at a relative gap of at least 0.20 and medium below; a later component makes it medium or hard by the same threshold; and an order that falls to the item key makes it hard. The 0.20 threshold is the smallest relative gap between adjacent evidence grades, $(1.0-0.8)/1.0$. The policy table names eleven distinct deciding-factor labels, of which the benchmark's label set uses nine: six commit labels (\code{deadline}, \code{evidence}, \code{impact}, \code{scope}, \code{explicit\_statement}, \code{least\_privilege}) and five decline labels (\code{insufficient\_separation}, \code{missing\_input}, \code{request\_ambiguous}, \code{stale\_explicit\_conflict}, \code{ask\_before\_crossing}). Appendix~\ref{app:policies} lists which policy produces each label; no record glosses the labels in prose, so that table is assembled from the policy table.

\subsection{Scenario families, perturbations and groups}
% source: families SC:45-71; axes SC:114-120; family size SC:162; perturbation kinds SC:277-289; comparable form SC:217-219; groups RP:63, LC:45-46.
Table~\ref{tab:families} lists the 18 scenario families. Each binds a task to one policy, draws fictional entities from a pool, and varies urgency (deadline pools per urgency class: critical 0, 1 or 2 days; high 1, 2, 3 or 5; medium 4, 7, 10 or 14; low 14, 21, 30 or 45), context length (0, 2 or 4 distractor items) and difficulty spacing (2, 1 or 0 days between adjacent deadlines). A family produces $n_{\text{base}} \times (1 + \sum \text{counts})$ examples: every base example is re-rendered under perturbations of six kinds, \code{paraphrase} (one of three alternative question wordings), \code{evidence\_order} and \code{context\_order} (reordered items), \code{irrelevant\_context} (up to two unused distractor sentences), \code{length} (one distractor added or removed) and \code{schema} (the prompt re-rendered as a Slack-style thread). Every perturbation must keep the decision's comparable form, the triple (ranking, deciding factor, abstained), equal to its base and must differ in content from every other example; generation fails otherwise. A family that reserves the JSON format cannot declare a \code{schema} perturbation, so no group straddles the format holdout. A base example and its perturbations share a group identifier, and the 1,350 examples form 294 such groups. Consistency across a group is therefore a property the data can test, and the group is the unit of dependence that the evaluation statistics must respect (Section~\ref{sec:protocol}).

% source: Kleos-Training-Data docs/scenarios.md, family table (SC:45-65) and the note on JSON-free families (SC:69-71).
\begin{table}[t]
\centering\scriptsize
\caption{The 18 scenario families of releases v0.0.6 and v0.0.7, with the examples each contributes to the release and to the test split, and the out-of-distribution shift it reserves. Task codes: npr notification prioritization, trt tool routing, mcb mission-control briefing, ctx context prioritization, rec recommendation generation, mcr memory-conflict resolution, wsr workspace reasoning. Three families have no JSON-format examples and so contribute nothing to the test split under the format holdout. Source: Kleos-Training-Data, \code{docs/scenarios.md}.}
\label{tab:families}
\begin{tabular}{lllrrl}
\toprule
Family & Task & Policy & Examples & Test & Reserved shift \\
\midrule
\code{ctx.budget\_pressure} & ctx & \code{rank\_by\_relevance\_over\_recency} & 120 & 0 & unseen entities \\
\code{ctx.relevance\_over\_recency} & ctx & \code{rank\_by\_relevance\_over\_recency} & 24 & 8 & reordered evidence \\
\code{mem.explicit\_over\_inferred} & mcr & \code{defer\_to\_explicit\_statement} & 100 & 35 & unseen formats \\
\code{mem.reliability\_over\_recency} & mcr & \code{rank\_by\_reliability\_over\_recency} & 88 & 32 & conflicting evidence \\
\code{mem.stale\_explicit\_conflict} & mcr & \code{defer\_to\_explicit\_statement} & 64 & 20 & conflicting evidence \\
\code{mem.unresolvable\_conflict} & mcr & \code{resolve\_or\_abstain\_on\_support} & 72 & 24 & conflicting evidence \\
\code{brief.across\_workspaces} & mcb & \code{respect\_workspace\_scope} & 100 & 35 & unseen formats \\
\code{brief.what\_needs\_a\_decision} & mcb & \code{rank\_by\_deadline\_then\_evidence} & 24 & 8 & context length shift \\
\code{notif.deadline\_vs\_evidence} & npr & \code{rank\_by\_deadline\_then\_evidence} & 24 & 8 & unseen entities \\
\code{notif.tied\_signals} & npr & \code{rank\_or\_abstain\_when\_close} & 100 & 30 & unseen formats \\
\code{rec.abstain\_without\_evidence} & rec & \code{rank\_or\_abstain\_when\_close} & 18 & 6 & unseen formats \\
\code{rec.verify\_before\_recommending} & rec & \code{verify\_when\_evidence\_weak} & 100 & 35 & unseen formats \\
\code{route.ask\_when\_underspecified} & trt & \code{ask\_when\_request\_ambiguous} & 90 & 30 & unseen formats \\
\code{route.least\_privilege} & trt & \code{prefer\_least\_privilege\_source} & 100 & 0 & unseen entities \\
\code{route.source\_matches\_question} & trt & \code{select\_by\_evidence\_need} & 88 & 32 & unseen formats \\
\code{wsp.absent\_in\_active} & wsr & \code{ask\_before\_crossing\_workspace} & 100 & 0 & unseen entities \\
\code{wsp.cross\_workspace\_request} & wsr & \code{respect\_workspace\_scope} & 120 & 40 & unseen formats \\
\code{wsp.scope\_boundary} & wsr & \code{respect\_workspace\_scope} & 18 & 6 & unseen formats \\
\midrule
Total & & & 1,350 & 349 & \\
\bottomrule
\end{tabular}
\end{table}

\subsection{Splitting and the held-out format}
% source: LC:45-72; SC:69-71, 295-297; configs/datasets/kleos_policy_v006.yaml:10-12, 40-42; SPEC:61-62; STG:110; EVAL:58-64; RM:157 (D2).
A group is never split across train, validation and test. Assignment is a stable hash: each group's rank is the first eight bytes of the SHA-256 of the split seed (42) and the group key, read as a fraction of $2^{64}$. Every release holds out the JSON input format: all 349 JSON-format examples go to the test split, and validation is drawn from the remaining formats, so early stopping never sees the out-of-distribution data the test split is meant to measure. The resulting split is 820 / 181 / 349 (recorded fractions 0.607 / 0.134 / 0.259). Three families (\code{ctx.budget\_pressure}, \code{route.least\_privilege} and \code{wsp.absent\_in\_active}) have no JSON-format examples and contribute nothing to the test split, so 15 families appear there, in 78 groups. A near-duplicate check at a Jaccard similarity of 0.85 or more reports 0 train--test near-duplicates for v0.0.6. The consequence for what the benchmark measures is recorded as deviation D2: test scenarios come from the same families as training, so the benchmark measures transfer across input format, not generalization to new scenarios.

\subsection{Releases and the reasoning traces}
% source: releases TDR:348-362; v0.0.6 changes TCL:154-172; abstain share RP:133-135; v0.0.7 V7:62-63, 93-132, 166-170 and TCL:113-122; test and benchmark hashes SPEC:54, EXP:477.
Releases are immutable and content-hashed. The two used here are v0.0.6 (2026-09-04; content hash \code{3cc9a744\ldots}) and v0.0.7 (2026-10-01; \code{b53afa42\ldots}), both of 1,350 examples in the same 820 / 181 / 349 split and both built from the same catalog with split seed 42. v0.0.6 introduced the relative close-call margin (replacing an absolute margin of 0.08), split the former \code{ask\_when\_underdetermined} policy into \code{verify\_when\_evidence\_weak} and \code{ask\_when\_request\_ambiguous}, and added the stale-statement rule; 338 of its 1,350 targets (25\%) abstain or ask. v0.0.7 keeps every example and decision of v0.0.6, and the identifiers of its JSON-format (test) examples, and adds two things to the training and validation answers only: a reasoning trace (schema 1.1) and a closing line ``What decided it: $\langle$label$\rangle$.'' on every answer, including declined ones. Each trace follows five steps and is capped at 1,200 characters: the rule the family's policy weighs; one line per item, in prompt order, mapping the prompt's wording to the policy's values (the design document's example reads ``Redpine: due in 2 days (0.33), reported (0.55), high (1.0) $\rightarrow$ 0.18''); each check the policy applied and its outcome; why the deciding factor decided it; and the conclusion, either the ranking or the failed check and the question to ask. All 1,001 training and validation answers carry one; the median trace is 560 characters long and the longest 878. The four decline labels that the test split expects thereby enter the training targets for the first time: \code{request\_ambiguous} (45 train / 15 validation), \code{stale\_explicit\_conflict} (32 / 12), \code{missing\_input} (25 / 5) and \code{insufficient\_separation} (22 / 4); the fifth, \code{ask\_before\_crossing}, enters in 100 training and validation answers combined, and no test item expects it. The test split of v0.0.7 is byte-identical to v0.0.6's (SHA-256 \code{a4decaaf\ldots}), and the benchmark derived from it (SHA-256 \code{a11ffad7\ldots}) is the same object for every run in this paper. Earlier releases, v0.0.1 to v0.0.5, record the pipeline's maturation; v0.0.5 removed 32 examples after a human review, described next.

\subsection{Quality and privacy record}
% source: detection layers TPE:20-25; gates STG:83-98, 144 and TDR:84-91; private-fact thresholds TPE:172-178; 32 checks LC:150; approvals PRV:91-94 and TDR:153-154; v0.0.5 review TCL:196-199; labelling quirk V7:67-69; synthetic PRV:42-44, 50.
The pipeline promotes examples through 14 gates, 13 of them mandatory (staging, schema and identifier integrity; four privacy gates; review presence and approval; provenance; corpus deduplication; evaluation leakage; and contract rendering; only the coverage gate can be bypassed). It applies 28 detection rules in three layers: 15 secret patterns that block, 11 personal-data patterns that redact and 2 structural heuristics that flag for review, beside a vault of registered entity literals. A private-fact classifier fails any example scoring 0.75 or more, routes scores from 0.4 to 0.75 to review and passes scores below 0.4. A sealed three-split release passes 32 verification checks. Two limits of this record matter for the paper. First, the approvals in every release to date were filled from the machine review's results; the only completed human review is v0.0.5's sample of 127 examples, of which 32 were rejected (explicit statements not honored, 3; memory conflicts decided on impact, 2; deadlines overridden, 4; ranking where the policy should have asked, 10; unjustified insufficient-separation abstentions, 13). Those categories drove the v0.0.6 policy changes, and no human review of v0.0.6 or v0.0.7 has been completed. Second, in 48 examples, 17 of them in the test split, the deciding factor is the largest raw gap between the top two items even when that gap favors the runner-up, a labelling quirk documented in the v0.0.7 design. No example in any release comes from a real user: every example was generated from the synthetic scenario catalog with fictional entities and captured through an offline mock backend, and the production-observation lane is rejected by the provenance gate.
